
Two factors are crossed:

\begin{itemize}
\item \textbf{Training type:} Cross-entropy (CE) fine-tuning vs. RL fine-tuning with execution rewards (REINFORCE)
\item \textbf{Inference mode:} Single-shot generation vs. K-step refinement (Self-Refine protocol, K=3)
\end{itemize}

This yields four conditions:

\begin{table}[htbp]
\centering
\resizebox{\columnwidth}{!}{%
\begin{tabular}{lll}
\toprule
Condition & Training & Inference \\
\midrule
CE-Single & Cross-entropy & Single-shot \\
CE-Refine & Cross-entropy & K=3 refinement \\
RL-Single & RL + execution & Single-shot \\
RL-Refine & RL + execution & K=3 refinement \\
\bottomrule
\end{tabular}
}
\end{table}

The interaction term---whether RL-Refine exceeds what CE-Refine and RL-Single would predict additively---directly tests the hypothesis of superadditivity.

\textbf{Statistical model:} A generalized linear mixed model (GLMM) with logit link:

$$\text{logit}(P(\text{pass})) = \beta_0 + \beta_T \cdot \text{Training} + \beta_R \cdot \text{Refinement} + \beta_{T \times R} \cdot \text{Training} \times \text{Refinement} + u_{\text{problem}}$$

where $u_{\text{problem}}$ is a random effect capturing problem-level difficulty. The success criterion is $\beta_{T \times R} > 0$ with $p < 0.05$ and odds ratio $\geq$ 1.2.
